%
% 6.006 problem set 4 solutions template
%
\documentclass[12pt,twoside]{article}

\input{macros-sp20}
\newcommand{\theproblemsetnum}{4}

\title{6.006 Problem Set \theproblemsetnum}

\begin{document}

\handout{Problem Set \theproblemsetnum}

\setlength{\parindent}{0pt}
\medskip\hrulefill\medskip

{\bf Name:} BUAA-TYZ

\medskip

{\bf Collaborators:} None

\medskip\hrulefill

%%%%%%%%%%%%%%%%%%%%%%%%%%%%%%%%%%%%%%%%%%%%%%%%%%%%%
% See below for common and useful latex constructs. %
%%%%%%%%%%%%%%%%%%%%%%%%%%%%%%%%%%%%%%%%%%%%%%%%%%%%%

% Some useful commands:
%$f(x) = \Theta(x)$
%$T(x, y) \leq \log(x) + 2^y + \binom{2n}{n}$
% {\tt code\_function}


% You can create unnumbered lists as follows:
%\begin{itemize}
%    \item First item in a list
%        \begin{itemize}
%            \item First item in a list
%                \begin{itemize}
%                    \item First item in a list
%                    \item Second item in a list
%                \end{itemize}
%            \item Second item in a list
%        \end{itemize}
%    \item Second item in a list
%\end{itemize}

% You can create numbered lists as follows:
%\begin{enumerate}
%    \item First item in a list
%    \item Second item in a list
%    \item Third item in a list
%\end{enumerate}

% You can write aligned equations as follows:
%\begin{align}
%    \begin{split}
%        (x+y)^3 &= (x+y)^2(x+y) \\
%                &= (x^2+2xy+y^2)(x+y) \\
%                &= (x^3+2x^2y+xy^2) + (x^2y+2xy^2+y^3) \\
%                &= x^3+3x^2y+3xy^2+y^3
%    \end{split}
%\end{align}

% You can create grids/matrices as follows:
%\begin{align}
%    A =
%    \begin{bmatrix}
%        A_{11} & A_{21} \\
%        A_{21} & A_{22}
%    \end{bmatrix}
%\end{align}

% You can include images and PDFs as follows:
% \includegraphics[width=0.5\textwidth]{img.jpg}

\begin{problems}

\problem  % Problem 1

\begin{problemparts}
\problempart % Problem 1a
Node(37): Scew = 2 - 0 = 2
Node(16): Scew = 1 - 3 = -2

\problempart % Problem 1b
\begin{center}
\begin{tikzpicture}
\tikzset{level distance=40pt, sibling distance=30pt}    
\Tree
[.\node[draw, circle]{$47$}; 
[.\node[draw, circle]{$16$}; 
[.\node[draw, circle]{$3$}; [.\node[draw, circle]{$2$}; ] ] 
[.\node[draw, circle]{$37$}; [.\node[draw, circle]{$28$}; ] ] 
] 
[.\node[draw, circle]{$85$}; 
[.\node[draw, circle]{$64$}; ] 
[.\node[draw, circle]{$86$}; [.\node[draw, circle]{$88$}; ] ] 
] 
]
\end{tikzpicture}
\end{center}

\problempart % Problem 1c
Just do a rotation.

\end{problemparts}

\newpage
\problem  % Problem 2

\begin{problemparts}
\problempart % Problem 2a
\begin{center}
    \begin{tikzpicture}
    \tikzset{level distance=40pt, sibling distance=30pt}    
    \Tree
    [.\node[draw, circle]{$4$}; 
    [.\node[draw, circle]{$12$}; 
    [.\node[draw, circle]{$21$}; ]
    [.\node[draw, circle]{$14$}; ]
    ] 
    [.\node[draw, circle]{$8$}; 
    [.\node[draw, circle]{$9$}; ] 
    [.\node[draw, circle]{$17$}; ] 
    ] 
    ]
    \end{tikzpicture}
\end{center}
Minheap
\problempart % Problem 2b
Maxheap
\problempart % Problem 2c
Swap (13, 0) and then (2, 0)
\problempart % Problem 2d
Minheap
\end{problemparts}

\problem  % Problem 3

\begin{problemparts}
\problempart % Problem 3a
\begin{itemize}
    \item Build a heap according to score: $O(|A|)$
    \item Delete max k times: $O(klog|A|)$
\end{itemize}


\problempart % Problem 3b
Use heap tree properity to search the heap.
\end{problemparts}

\problem  % Problem 4
$s_i$: address of farm $c_i$: capacity $b_j$: name of building $d_j$: demand of building
\begin{itemize}
    \item A hash table to record $(s_i, node)$ pair.
    \item For each node, a hash table to record $(b_j, "")$ pair.
    \item An max heap to store all node based on $c_i$.
\end{itemize}
Build: $O(n)$
Power on: Find a node corrsponding to the condition $O(logn)$, add the name to 
the hash table of the node. Change the capacity to $c_j - d_j$
Power off: Find the node $O(logn)$, delete it in the hash table: $O(1)_{am}$
Customers: Use hash table to find the node and return the hash table of the node.

\newpage
\problem  % Problem 5
Pass problem 5 and 6, use the time to see the implementation of AVL Tree.

\problem  % Problem 6
\begin{problemparts}
\problempart % Problem 6a
\problempart % Problem 6b
\problempart % Problem 6c
\problempart Submit your implementation to {\small\url{alg.mit.edu}}.
\end{problemparts}

\end{problems}

\end{document}
